% Generated draft: requires mathematical and source review.
\documentclass[11pt]{article}
\usepackage[margin=25mm]{geometry}
\usepackage{fontspec}
\setmainfont{Latin Modern Roman}
\usepackage{amsmath,amssymb,mathtools,unicode-math}
\setmathfont{Latin Modern Math}
\usepackage{enumitem,needspace,xurl,hyperref}
\setlength{\parindent}{0pt}
\setlength{\parskip}{5pt}
\setlength{\emergencystretch}{3em}
\newcommand{\N}{\mathbb{N}}
\newcommand{\Z}{\mathbb{Z}}
\newcommand{\Q}{\mathbb{Q}}
\newcommand{\R}{\mathbb{R}}
\newcommand{\C}{\mathbb{C}}
\newcommand{\F}{\mathbb{F}}
\newcommand{\A}{\mathbb{A}}
\newcommand{\PP}{\mathbb{P}}
\newcommand{\E}{\mathbb{E}}
\newcommand{\eps}{\varepsilon}
\newcommand{\dd}{\,\mathrm{d}}
\begin{document}
\section*{UnsolvedMath ID: 30004777}
\textbf{Source problem code:} catalogue entry
\subsection{Definitions and mathematical statement}
The symbol \(H\) denotes the group in the supplied assertion. A formal degree is the scalar in the Schur orthogonality relation for an appropriate square-integrable representation, with the Haar measure and any convention concerning the center fixed. A Plancherel density describes a Plancherel measure relative to a specified reference measure. Neither the required normalizations nor the precise representation and parameter spaces are supplied here.

The extracted assertion is: for every \(H\) that is a symplectic group or an even special orthogonal group and satisfies the hypotheses of the referenced local Langlands framework, the formal-degree and Plancherel-density conjectures hold for \(H\).

This is an incomplete assertion schema, rather than a fully specified mathematical statement: the record does not provide the base field, allowable forms and ranks of \(H\), local Langlands hypotheses, relevant representations and parameters, or either conjectural formula and its normalizations. Those missing hypotheses cannot be reconstructed from the supplied text.

The supplied verification says that the report presents this assertion as \emph{Theorem 1.3}, a proved special case, rather than an open problem; the record marks its statement status as \texttt{unrecoverable}. No separate unresolved mathematical target is specified. The supplied broader assessment says that these cases, alongside unitary, odd orthogonal, regular, and unipotent cases, do not settle the general formal-degree conjecture. These are supplied source characterizations, not independently verified mathematical conclusions.
\subsection{Short English statement}
The record attributes the formal-degree and Plancherel-density assertions for symplectic and even special orthogonal groups to a proved theorem under an unspecified local Langlands framework. Its missing hypotheses and formulas prevent a precise standalone statement. The supplied partially solved label concerns a broader conjectural context whose unresolved target is not specified.
\subsection{Sources}
\begin{itemize}
\item Local UnsolvedMath registry, permanent ID \(30004777\), exact source code \texttt{catalogue entry}, title \emph{Formal Degrees for Symplectic and Even Orthogonal Groups}. This draft uses the supplied complete source record. Its proposed year is 2021; no named proposer is supplied.
\item Supplied report citation: \emph{Automorphic Forms, Geometry and Arithmetic} (2022), Oberwolfach Reports, pp.\ 2089--2156, DOI 10.4171/owr/2021/39. \url{https://doi.org/10.4171/owr/2021/39}. The supplied verification identifies the extracted assertion as Theorem 1.3 and describes it as a proved special case. This identification has not been independently checked.
\item Supplied literature source: \emph{Oberwolfach Report 39/2021}. \url{https://doi.org/10.4171/owr/2021/39}.
\item Supplied literature source: \emph{Hiraga--Ichino--Ikeda conjecture for classical groups}. \url{https://arxiv.org/abs/2508.08470}. No authors or specific theorem references for this item are supplied.
\item Literature assessment dated exactly \texttt{2026-08-22}: the supplied status is \texttt{partially\_solved}. The supplied assessment describes symplectic and even special orthogonal groups as established cases, alongside unitary, odd orthogonal, regular, and unipotent cases, while saying that the general formal-degree conjecture remains unsettled. This is supplied dated curation derived from the Oberwolfach material, not independently verified literature assessment; its time-sensitive status and cited claims require independent checking.
\end{itemize}
\end{document}
